\documentclass[a4paper,10pt]{article}
\usepackage{amsmath, amsfonts, amssymb, physics}
\usepackage{revsymb}

\usepackage{amsmath,amsfonts,amsthm,physics,braket,indentfirst}
\usepackage{float}

\usepackage[a4paper,
            left=3cm,
            right=3cm,
            top=3cm,
            bottom=3cm]{geometry}

\usepackage{graphicx}

\usepackage{hyperref}
\hypersetup{
	colorlinks=true,       	
	linkcolor=blue,          	
	citecolor=blue,        		
	urlcolor=blue}
\usepackage{url}

\newtheorem{alg}{Algorithm}
\newtheorem{exe}{Exercise}

\title{Synchronization}
\author{M. Gil de Oliveira}
\date{\today}

\begin{document}

\maketitle

\section{Introduction}

Consider the following model for the synchronization between two separated parties $j = 1,2$: each site has a clock which ticks at a rate $f_j$. They believe their clocks are very close to a certain frequency $f_0$ (in our case, $f_0 = 10\text{MHz}$), but, actually, there is a frequency difference $\Delta f = f_2 - f_1$ which is tiny but not negligible.

For an event that happens at a certain time $t$, each party does not measure $t$ directly, but the amount of clock ticks
\begin{equation}
    N_j = f_j t \approx f_0 t
\end{equation}
Imposing that $t$ is the same at both sites leads us to the relation
\begin{equation}
    \frac{N_1}{f_1} = \frac{N_2}{f_2}
\end{equation}

Defining the delay $\tau = (N_2 - N_1) / f_0$, we may write
\begin{equation}
    \tau = \frac{\Delta f}{f_1}\frac{N_1}{f_0} = \frac{\Delta f}{f_0} t.
\end{equation}
More generally, when $\Delta f(t)$ is a function of $t$, we can express this rule as either
\begin{equation}
    \tau (t) = \int_{0}^{t} \frac{\Delta f(s)}{f_0} ds
\end{equation}
or
\begin{equation}
    \frac{d \tau}{dt} =  \frac{\Delta f(t)}{f_0}.
\end{equation}
In either case, we see that $\frac{\Delta f(t)}{f_0}$ is the quantity governing the clock drift.

\section{Experiment}

\end{document}
